% SIAM Shared Information Template
% This is information that is shared between the main document and any
% supplement. If no supplement is required, then this information can
% be included directly in the main document.


% Packages and macros go here
\usepackage{lipsum}
\usepackage{amsfonts}
\usepackage{graphicx}
\usepackage{epstopdf}
\usepackage{algorithm}      % for floating algorithm environment
\usepackage{algpseudocode}  % for the algorithmicx syntax (\State, \While, etc.)
\usepackage{mathtools}
\usepackage{booktabs}
\usepackage{subcaption}
\usepackage[scr=rsfs]{mathalpha}
\usepackage{multirow}
\usepackage{bbm}
\usepackage{longtable,booktabs}
\ifpdf
  \DeclareGraphicsExtensions{.eps,.pdf,.png,.jpg}
\else
  \DeclareGraphicsExtensions{.eps}
\fi

% Prevent itemized lists from running into the left margin inside theorems and proofs
\usepackage{outlines}
\usepackage{enumitem}
\setlist[enumerate]{leftmargin=.5in}
\setlist[itemize]{leftmargin=.5in}

% Add a serial/Oxford comma by default.
\newcommand{\creflastconjunction}{, and~}

% Used for creating new theorem and remark environments
\newsiamremark{assumption}{Assumption}
\newsiamremark{remark}{Remark}
\newsiamremark{hypothesis}{Hypothesis}
\crefname{hypothesis}{Hypothesis}{Hypotheses}
\newsiamthm{claim}{Claim}

% Sets running headers as well as PDF title and authors
\headers{Distribution-Free Robust Predict-Then-Optimize in Function Spaces}{Y. Patel and A. Tewari}

% Title. If the supplement option is on, then "Supplementary Material"
% is automatically inserted before the title.
\title{Distribution-Free Robust Predict-Then-Optimize in Function Spaces\thanks{Submitted to the editors Feb 8, 2026.
\funding{We acknowledge the support of NSF via grant DMS-2413089.}}}

% Authors: full names plus addresses.
\author{Yash Patel\thanks{University of Michigan, Ann Arbor, MI 
  (\email{yppatel@umich.edu}).}
\and Ambuj Tewari\thanks{University of Michigan, Ann Arbor, MI 
  (\email{tewaria@umich.edu}).}}

\usepackage{subcaption} % in your preamble

\usepackage{amsopn}
\DeclareMathOperator{\diag}{diag}
\DeclareMathOperator*{\argmax}{arg\,max}
\DeclareMathOperator*{\argmin}{arg\,min}
\DeclareMathOperator{\logit}{logit}
\DeclarePairedDelimiter{\ceil}{\lceil}{\rceil}
\DeclarePairedDelimiter\floor{\lfloor}{\rfloor}
\newcommand{\indep}{\perp \!\!\! \perp}

% User-defined shortcuts
\newcommand{\rtwo}{\mathbb{R}^2}
\newcommand{\rone}{\mathbb{R}}
\newcommand{\qone}{\mathbb{Q}}
\newcommand{\qtwo}{\mathbb{Q}^2}
\newcommand{\nat}{\mathbb{N}}
\newcommand{\ex}{\mathbb{E}}
\newcommand{\rt}{\rightarrow}
\newcommand{\lt}{\leftarrow}
\newcommand{\cW}{\mathcal{W}}
\newcommand{\cU}{\mathcal{U}}
\newcommand{\cE}{\mathcal{E}}
\newcommand{\cC}{\mathcal{C}}
\newcommand{\T}{\top}
\newcommand{\normal}{\mathcal{N}}
\newcommand{\vor}{\mathcal{V}}
\newcommand{\st}{\mid}
\newcommand{\abs}[1]{\left|#1\right|}
\newcommand{\set}[1]{\left\{#1\right\}}
\newcommand{\norm}[1]{\left\lVert #1\right\rVert}
\newcommand{\obar}[1]{\overline{#1}}
\newcommand{\cov}{\text{cov}}

\renewcommand{\algorithmicrequire}{\textbf{Input:}}
\renewcommand{\algorithmicensure}{\textbf{Output:}}

\crefname{subsection}{Section}{Sections}
\Crefname{subsection}{Section}{Sections}
\crefname{subsubsection}{Section}{Sections}
\Crefname{subsubsection}{Section}{Sections}
\crefname{assumption}{Assumption}{Assumptions}
\Crefname{assumption}{Assumption}{Assumptions}

%%%% HELPER CODE FOR DEALING WITH EXTERNAL REFERENCES ON OVERLEAF
% (from an answer by cyberSingularity at http://tex.stackexchange.com/a/69832/226)
%%%
\makeatletter
\newcommand*{\addFileDependency}[1]{% argument=file name and extension
  \typeout{(#1)}% latexmk will find this if $recorder=0 (however, in that case, it will ignore #1 if it is a .aux or .pdf file etc and it exists! if it doesn't exist, it will appear in the list of dependents regardless)
  \@addtofilelist{#1}% if you want it to appear in \listfiles, not really necessary and latexmk doesn't use this
  \IfFileExists{#1}{}{\typeout{No file #1.}}% latexmk will find this message if #1 doesn't exist (yet)
}
\makeatother

\newcommand*{\myexternaldocument}[1]{%
    \externaldocument{#1}%
    \addFileDependency{#1.tex}%
    \addFileDependency{#1.aux}%
}
%%% END HELPER CODE

%%% Local Variables: 
%%% mode:latex
%%% TeX-master: "ex_article"
%%% End: 
